% !TEX root = micro_Lie_theory.tex
% source: [local source path]
% converted_by: Codex
% date: 2026-06-01
% compile_status: failed (xelatex: hyperref wrong driver `pdftex` in module preamble)
% known_issues: xelatex module fails at micro_Lie_theory.tex with Wrong driver option `pdftex`

\begin{abstract}

Група Лі — старий абстрактний математичний об'єкт, що походить з XIX століття, коли математик Софіса Лі заклав основи теорії груп неперервних перетворень. 
Його вплив поширився на різні галузі науки й технологій значно пізніше.
У роботі протягом останніх років спостерігається важливий тренд зростання її використання, принаймні в областях оцінювання, особливо в оцінюванні руху для навігації.
Водночас для переважної більшості фахівців з робототехніки групи Лі залишаються надто абстрактними побудовами, а отже важкими для розуміння й практичного застосування.

У задачах оцінювання для робототехніки часто немає потреби повністю використовувати всю силу теорії, тому потрібен добір матеріалу.
У цій статті ми пройдемо через основні принципи теорії Лі з метою подати чіткі й корисні ідеї, залишивши в стороні значний корпус додаткової теорії.
Навіть після такого скорочення поданий матеріал довів надзвичайну ефективність у сучасних алгоритмах оцінювання для робототехніки, особливо в задачах SLAM, візуальної одометрії та подібних підходів.

Поряд із цією мікро-теорією Лі ми подаємо розділ із кількома прикладами застосувань та великий довідник формул для основних груп Лі, що використовуються в робототехніці, включно з більшістю матриць Якобі та способів їх зручного маніпулювання.
Ми також представляємо нову C++ бібліотеку на чистих шаблонах, яка реалізує весь функціонал, описаний у цьому матеріалі.


\end{abstract}


